\documentclass[../main.tex]{subfiles}

\begin{document}

\ifSubfilesClassLoaded{

    %\tableofcontents
    
    %\setcounter{chapter}{}
}{}

\chapter{ Baire Space }

\begin{definition}{}{}
A space $X$ is said to be a \dfntxt{Baire space} if the following condition holds:
Given any countable collection $\{A_n\}$ of closed sets of $X$ each of which has empty interior in $X$, their union $\bigcup A_n$ also has empty interior in $X$.
\end{definition}

\begin{definition}{}{}
  A subset $A$ of a space $X$ was said to be of the \dfntxt{first category} in $X$ if it was contained in the union of a countable collection of closed sets of $X$ having empty interiors in $X$; otherwise, it was said to be of the \dfntxt{second category} in $X$. 
\end{definition}

\begin{theorem}{}{}
    A space $X$ is a \dfntxt{Baire space} if and only if every nonempty open set in $X$ is of the \dfntxt{second category}.
\end{theorem}
\end{document}